\documentclass[12pt]{article}
\usepackage{amsmath,amssymb,amsfonts}
\usepackage[margin=1in]{geometry}

\begin{document}

\section*{Problem 5.1}

\textbf{Problem:} Prove that $\displaystyle\int_a^b f(x)\,\delta(x)\,dx = 0$ for $a \neq 0$ in two ways.

\bigskip
\textbf{Solution:}

We need to show that $\int_a^b f(x)\,\delta(x)\,dx = 0$ when $0 \notin [a,b]$, i.e., when $a$ and $b$ are both positive or both negative (or more precisely, when the origin is not in the interval of integration).

\subsection*{Method 1: Using the defining property of the delta function}

The Dirac delta function is defined by its action on test functions:
\[
\int_{-\infty}^{\infty} f(x)\,\delta(x)\,dx = f(0).
\]
More generally, for any interval $[a,b]$:
\[
\int_a^b f(x)\,\delta(x)\,dx = \begin{cases} f(0) & \text{if } a < 0 < b, \\ 0 & \text{if } 0 \notin [a,b]. \end{cases}
\]

If $a \neq 0$ and the interval $[a,b]$ does not contain the origin (for example, $0 < a < b$ or $a < b < 0$), then the delta function is zero throughout the entire interval of integration. Since $\delta(x) = 0$ for all $x \neq 0$, and $x = 0$ is not in $[a,b]$, we have
\[
\int_a^b f(x)\,\delta(x)\,dx = 0.
\]

\subsection*{Method 2: Using the delta function as a limit of a sequence}

Consider the delta function as the limit of the sequence of functions:
\[
\delta_n(x) = \begin{cases} n/2 & \text{if } |x| < 1/n, \\ 0 & \text{if } |x| \geq 1/n. \end{cases}
\]

For an interval $[a,b]$ with $0 \notin [a,b]$, there exists a minimum distance $d = \min(|a|, |b|) > 0$ from the origin to the interval. For all $n > 1/d$, the support of $\delta_n(x)$, namely $(-1/n, 1/n)$, does not intersect $[a,b]$. Therefore $\delta_n(x) = 0$ for all $x \in [a,b]$ when $n > 1/d$, and
\[
\int_a^b f(x)\,\delta_n(x)\,dx = 0 \quad \text{for all } n > 1/d.
\]

Taking the limit:
\[
\int_a^b f(x)\,\delta(x)\,dx = \lim_{n \to \infty} \int_a^b f(x)\,\delta_n(x)\,dx = 0.
\]

\end{document}
